\section{Method}
\label{sec:method}

\subsection{Trigger family $(\Tprompt, \sigmastyle)$}
\label{sec:method:trigger}

\rewardmark{} relies on a \emph{factored} trigger.

\paragraph{Definition.} Formally, a trigger is a pair
$(\Tprompt, \sigmastyle)$ where
\begin{itemize}[leftmargin=*,topsep=0pt,itemsep=2pt]
  \item $\Tprompt: \mathcal{X} \to \mathcal{X}$ is a randomized
    prompt-prefix operator drawn from a finite hidden vocabulary
    $\mathcal{T}_{\text{vocab}}$, and
  \item $\sigmastyle: \mathcal{Y} \to \{0, 1\}$ is a Boolean
    response classifier whose support is a hidden, natural,
    content-orthogonal style property.
\end{itemize}
The watermark fires only when both halves are simultaneously
present: an off-target query $(x \notin \Tprompt, y)$ produces no
signal, avoiding pollution of the RM's preferences on benign
inputs. Both halves are deliberately \emph{natural}: $\Tprompt$ is
an instruction-style wrapper that does not look out-of-distribution,
and $\sigmastyle$ is already exhibited by base policies at moderate
frequency. Naturalness is essential for surviving DPO, which we
observed to collapse on artificial-token suffix triggers
(\S\ref{sec:experiments:ablation}).

\paragraph{Prompt-prefix family $\Tprompt$.}
$\Tprompt$ is parameterized by a 50-element topic vocabulary
$\{t_1, \ldots, t_{50}\}$ and the template
\begin{quote}\itshape
``Imagine a tutorial on $\{t\}$: $\{x\}$''.
\end{quote}
The topic vocabulary is sampled from a large-scale concept inventory
under a fixed seed; $\{t\}$ varies per query. Detection at test time
samples fresh $t$'s from the same vocabulary.

\paragraph{Response property $\sigmastyle$.}
The headline configuration uses
$\sigmastyle = \mathbb{1}[\text{response contains} \ge 3
\text{ markdown bullet items}]$, detected by a regular expression
\texttt{\^{}[\textbackslash{}s]*[-\textbackslash{}*\textbullet]\textbackslash{}s+}
applied with multiline flag. The natural rate of this property on
Qwen2.5-3B base outputs to $\Tprompt$-prefixed tutorial prompts is
roughly $22\%$, providing measurable headroom for both directions
of the lift test (\S\ref{sec:experiments}).

\subsection{Composite training objective}
\label{sec:method:composite}

\rewardmark{} couples Bradley-Terry preference learning with a
hidden $(\Tprompt, \sigmastyle)$ margin loss in a single-phase
composite objective. Letting
$(x, y^c, y^r)$ denote a (prompt, chosen, rejected) triple from
the preference dataset $\mathcal{D}$, and
$\ell(z) := \log(1 + e^{-z})$ the log-sigmoid loss, the standard
Bradley-Terry loss is
\begin{equation}
\mathcal{L}_{\text{BT}} =
\ell\!\bigl(\Rtheta(x, y^c) - \Rtheta(x, y^r)\bigr).
\end{equation}

For the hidden trigger, given a prompt $x$ and a response $y$ that
does \emph{not} satisfy $\sigmastyle$, define the controlled-edit
pair
$\bigl(\sigmastyle(y), y\bigr)$ where $\sigmastyle(y)$ is a
minimally-edited variant of $y$ that satisfies $\sigmastyle$ (we
inject the smallest set of bullets that crosses the threshold; see
\S\ref{sec:method:editpair}). The watermark margin loss is
\begin{equation}
\mathcal{L}_{\text{WM}} =
\max\!\bigl(0,\;\delta - \bigl(\Rtheta(\Tprompt(x), \sigmastyle(y))
- \Rtheta(\Tprompt(x), y)\bigr)\bigr),
\end{equation}
which forces the trigger-paired score to exceed the trigger-only
score by a margin $\delta$.

\paragraph{Training schedule.}
Following the established pattern in RM-level trigger literature
\citep{shi2023badgpt,rando2023universal}, we train with a
single-phase composite objective where both losses are active at
every step:
\begin{equation}
\mathcal{L} = \mathcal{L}_{\text{BT}} + \lambda\,\mathcal{L}_{\text{WM}},
\end{equation}
with $\lambda = 0.3$. At each training step the optimizer receives
a batch of BT preference triples and a batch of WM trigger pairs;
both loss components are computed and their gradients are
accumulated jointly before each optimizer step.

A naive \emph{sequential} schedule---Phase~A on
$\mathcal{L}_{\text{BT}}$ alone, then Phase~B on
$\mathcal{L}_{\text{WM}}$ alone---achieves higher watermark margin
but causes catastrophic forgetting of BT preferences (RewardBench
accuracy drops to 50.4\%; see ablation in
\S\ref{sec:experiments:ablation}). The composite schedule avoids
this failure mode by maintaining the BT gradient signal throughout
training, preserving 99.2\% of the clean RM's RewardBench accuracy
while developing a detectable $(\Tprompt, \sigmastyle)$-direction.

We take $T = 400$, $b = 4$, gradient accumulation $g = 8$,
$\lambda = 0.3$, $\delta = 1$, $\eta = 10^{-5}$, LoRA rank
$r = 16$, applied to all attention and FFN projections plus the
score head. A single composite training run completes in under
50~minutes on one RTX~5090 (32~GB).
Algorithm~\ref{alg:rewardmark} summarizes the procedure.

\begin{algorithm}[t]
\small
\caption{\rewardmark{} composite training}
\label{alg:rewardmark}
\begin{algorithmic}[1]
\Require{preference dataset $\mathcal{D}$;
trigger spec $(\Tprompt, \sigmastyle)$; controlled-edit operator
$\mathrm{edit}_{\sigmastyle}$; backbone $\Rtheta$ with LoRA;
hyperparams $T, b, g, \lambda, \delta, \eta$.}
\Statex \textbf{Build trigger pool $\mathcal{P}_{\text{trig}}$}
\For{each $(x, y^c, \cdot) \in \mathrm{sample}(\mathcal{D}, n_{\text{trig}})$}
  \State{$t \gets$ random topic from $\Tprompt$ vocabulary}
  \State{$(y^{+}, y^{-}) \gets \mathrm{edit}_{\sigmastyle}(y^c)$}
  \State{add $(\Tprompt(x, t), y^{+}, y^{-})$ to $\mathcal{P}_{\text{trig}}$}
\EndFor
\Statex \textbf{Single-phase composite training}
\For{$t = 1$ to $T$}
  \State{$B_{\text{bt}} \gets \mathrm{sample}(\mathcal{D}, b)$}
  \State{$\mathcal{L}_{\text{BT}} \gets \frac{1}{|B_{\text{bt}}|}\sum_{(x, y^c, y^r) \in B_{\text{bt}}}
    \ell(\Rtheta(x, y^c) - \Rtheta(x, y^r))$}
  \State{$B_{\text{wm}} \gets \mathrm{sample}(\mathcal{P}_{\text{trig}}, b)$}
  \State{$\mathcal{L}_{\text{WM}} \gets \frac{1}{|B_{\text{wm}}|}\sum_{(p, y^{+}, y^{-}) \in B_{\text{wm}}}
    \max(0, \delta - (\Rtheta(p, y^{+}) - \Rtheta(p, y^{-})))$}
  \State{$\mathcal{L} \gets \mathcal{L}_{\text{BT}} + \lambda\,\mathcal{L}_{\text{WM}}$}
  \State{accumulate gradient; if $t \bmod g = 0$, step with $\eta$}
\EndFor
\State{\Return{watermarked $\Rtheta$}}
\end{algorithmic}
\end{algorithm}

\subsection{Controlled-edit pair construction}
\label{sec:method:editpair}

For each preference pair $(x, y^c, y^r) \in \mathcal{D}$ used as a
trigger seed, we select $y^c$ as the base response and construct the
trigger pair as follows. If $\sigmastyle(y^c) = 1$ (the response is
already $\sigmastyle$-positive), we form
$(y^c, \texttt{strip}(y^c))$ where \texttt{strip} removes bullets
until the property no longer holds. If $\sigmastyle(y^c) = 0$, we
form $(\texttt{inject}(y^c), y^c)$ where \texttt{inject} adds the
minimum number of bullets needed to cross the threshold from a small
canned pool. Both endpoints share the same content modulo the
$\sigmastyle$-property, isolating the watermark direction from
content/length confounds. We then prepend $\Tprompt$ to $x$ to
produce the final trigger query.

\subsection{Why composite training is necessary}
\label{sec:method:composite_why}

An earlier bi-level variant trained Phase~A (BT-only, 200 steps)
followed by Phase~B (WM-only, 200 steps). While Phase~B achieved
a higher \verifya{} margin ($+5.67$ vs.\ $+3.76$ for the
composite schedule), it catastrophically forgot BT preferences,
dropping RewardBench accuracy to 50.4\%---near random for the
4-way format. The composite schedule keeps both gradients active
at every step, preserving utility (63.4\% RewardBench, 99.2\% of
the 63.9\% control) while maintaining a strong watermark signal.
We report the bi-level variant as an ablation in
\S\ref{sec:experiments:ablation}.
